\section{Entendimiento del Negocio}
El análisis de la calidad educativa en Colombia no puede abordarse de manera aislada, ya que está condicionado por múltiples factores estructurales, entre ellos las condiciones socioeconómicas, la conectividad digital y las características territoriales. En este sentido, comprender el problema educativo implica analizar cómo estas variables interactúan y generan brechas persistentes en el desempeño académico de los estudiantes, especialmente en los municipios con mayores niveles de vulnerabilidad.
\subsection{Contexto territorial y estructural}
La categorización de los municipios en Colombia responde a criterios relacionados con población, capacidad fiscal e ingresos corrientes, lo que permite identificar niveles diferenciados de desarrollo territorial. En este contexto, los municipios de categorías 4, 5 y 6 se caracterizan por presentar menor capacidad institucional, alta dependencia de transferencias del Estado y limitaciones en infraestructura, lo que los ubica dentro de los territorios con mayores niveles de vulnerabilidad (Departamento Nacional de Planeación, 2025).
Estos municipios, en su mayoría rurales o dispersos, enfrentan barreras significativas en términos de acceso a servicios básicos y conectividad digital. Por ejemplo, la penetración de internet en zonas rurales es considerablemente menor que en zonas urbanas, lo que limita el acceso a herramientas educativas digitales y reduce las oportunidades de aprendizaje autónomo (Fundación Empresarios por la Educación [ExE], 2025).
Además, la dispersión geográfica y las dificultades de acceso físico incrementan los costos de implementación de infraestructura tecnológica, lo que dificulta la ejecución efectiva de programas nacionales de conectividad y educación digital (Bancolombia, s.f.).
\subsection{Indicadores macroeconómicos y su impacto}
El comportamiento de la economía nacional condiciona la capacidad del Estado para invertir en educación y conectividad, especialmente en territorios con baja generación de ingresos propios (Departamento Administrativo Nacional de Estadística [DANE], 2024).
Para el año 2024, el crecimiento del Producto Interno Bruto (PIB) en Colombia fue del 1,6\%, con una proyección de recuperación hacia el 2,6\% en 2025, lo cual refleja una dinámica económica moderada que limita la expansión acelerada de la inversión pública en infraestructura educativa
Por otro lado, la inflación alcanzó niveles cercanos al 7,16\%, afectando el costo de dispositivos tecnológicos y servicios asociados a la conectividad, lo que impacta directamente la capacidad de los hogares para acceder a herramientas digitales necesarias para el aprendizaje.
El desempleo, proyectado en 8,3\%, presenta niveles aún más altos en zonas rurales, donde la informalidad laboral supera el 85\%, generando inestabilidad en los ingresos familiares y afectando la permanencia escolar de los estudiantes.
Además, el nivel de deuda pública, cercano al 71,5\% del PIB, restringe el margen fiscal para inversiones directas en infraestructura de conectividad en zonas apartadas, lo que refuerza las desigualdades territoriales.
Estos indicadores evidencian que el problema educativo no puede analizarse de forma aislada, sino como parte de un sistema económico que condiciona la disponibilidad de recursos y la capacidad de implementación de políticas públicas.
\subsection{Pobreza multidimensional y brechas sociales}
El Índice de Pobreza Multidimensional (IPM) permite comprender las múltiples privaciones que afectan a los hogares colombianos. En 2024, la incidencia nacional se ubicó en 11,5\%, sin embargo, en zonas rurales alcanza valores cercanos al 24,3\%, lo que evidencia una brecha significativa frente a las cabeceras municipales (Departamento Administrativo Nacional de Estadística [DANE], 2024).
Dentro de las dimensiones del IPM, el bajo logro educativo afecta aproximadamente al 68\% de los hogares rurales, lo que indica que los estudiantes provienen de entornos con limitadas capacidades de acompañamiento académico. Asimismo, factores como el trabajo informal, el retraso escolar y las barreras de acceso a servicios de salud influyen indirectamente en el desempeño educativo, ya que afectan tanto la estabilidad del hogar como la continuidad en los procesos de formación.
Estas condiciones configuran un entorno en el cual la desigualdad no solo es económica, sino también educativa, reproduciendo brechas intergeneracionales en el acceso a oportunidades.
\subsection{Calidad educativa y resultados Saber 11}
Los resultados de las pruebas Saber 11 evidencian las desigualdades existentes en el sistema educativo colombiano. Para el año 2024, el puntaje global promedio fue de 256 puntos, mostrando una leve mejora frente al año anterior; sin embargo, esta tendencia no implica una reducción real de las brechas (Fundación Empresarios por la Educación [ExE], 2025).
La diferencia entre zonas urbanas y rurales alcanza aproximadamente 26,1 puntos, lo que representa una de las brechas más altas de los últimos años.
También, el análisis por tipo de institución muestra que los establecimientos oficiales rurales presentan puntajes significativamente inferiores en comparación con instituciones urbanas o privadas, lo que refleja limitaciones estructurales en la calidad del sistema educativo en estos territorios.
Los datos del MEN confirman estas desigualdades a nivel de proceso educativo. La cobertura neta en primaria se mantiene alrededor del 90\% a nivel nacional, mientras que en educación media cae a valores de entre 40\% y 45\%, lo que evidencia una brecha progresiva de acceso conforme avanza el nivel escolar. En términos de permanencia, la tasa de deserción promedio nacional oscila entre el 3\% y 4\%, siendo la educación media el nivel más crítico con tasas superiores al 5\% en promedio y casos extremos que superan el 10\% en municipios rurales. Adicionalmente, los indicadores muestran diferencias de hasta 30 puntos porcentuales en cobertura neta entre departamentos desarrollados y periféricos, lo que complementa las brechas observadas en los resultados Saber 11.
En departamentos con altos niveles de vulnerabilidad, como Chocó, los puntajes promedio pueden descender hasta valores cercanos a 210 puntos, mientras que en ciudades como Bogotá superan los 270 puntos, evidenciando profundas desigualdades territoriales
\subsection{Brecha digital y conectividad}
La brecha digital en Colombia es un fenómeno multidimensional que no se limita al acceso a infraestructura, sino que incluye factores como habilidades digitales, motivación y uso efectivo de las tecnologías.
El Índice de Brecha Digital muestra que, aunque ha habido avances en el acceso material, persisten dificultades en el desarrollo de habilidades digitales, siendo esta una de las dimensiones más críticas en los últimos años (Ministerio de Tecnologías de la Información y las Comunicaciones [MinTIC], 2023).
En zonas rurales, la conectividad es limitada e inestable, con una proporción significativa de usuarios que califican el servicio como deficiente, lo que dificulta el uso de plataformas educativas, simulaciones y recursos digitales necesarios para el aprendizaje.
Además, la existencia de más de 12 millones de personas sin acceso a internet refleja la magnitud de la brecha digital en el país, especialmente concentrada en territorios rurales y de baja densidad poblacional.
El análisis de los datos del MEN sobre sedes educativas oficiales muestra que el porcentaje de sedes conectadas a internet ha tenido una evolución positiva durante la última década, pasando de niveles inferiores al 30\% en 2011 a superar el 60\% en años recientes a nivel nacional. Sin embargo, la heterogeneidad territorial persiste: el promedio nacional de sedes conectadas es del 34.8\% (mediana: 27.3\%), y varios departamentos se mantienen por debajo del 40\% de conectividad. Adicionalmente, el análisis de correlación entre conectividad y deserción escolar arroja un valor de -0.049, sugiriendo que a mayor porcentaje de sedes conectadas, tiende a observarse menor deserción, aunque la magnitud de la correlación indica que la conectividad por sí sola no explica el fenómeno en su totalidad. La correlación entre conectividad y cobertura neta es de +0.144, coherente con la hipótesis de que la infraestructura tecnológica actúa como factor protector del acceso educativo.
\subsection{Impacto post-pandemia}
La pandemia de COVID-19 profundizó las desigualdades existentes en el sistema educativo. El cierre de instituciones educativas y la transición a modelos de educación remota afectaron de manera desproporcionada a los estudiantes con menor acceso a conectividad (Instituto Colombiano para la Evaluación de la Educación [ICFES], 2025).
Estudios han mostrado que los estudiantes con acceso a computador e internet obtuvieron puntajes significativamente superiores en comparación con aquellos que no contaban con estos recursos, evidenciando el papel clave de la tecnología en los procesos de aprendizaje
Asimismo, la pérdida de aprendizaje asociada a la reducción de clases presenciales ha tenido efectos persistentes, especialmente en contextos rurales, donde las condiciones para la educación remota fueron limitadas.
\subsection{Definición del problema de negocio}
En este contexto, el problema de negocio se puede definir como la dificultad para traducir el acceso a conectividad en mejoras reales en el desempeño educativo en municipios con altos niveles de vulnerabilidad.
Es decir, aunque se han realizado inversiones en infraestructura digital, estas no siempre se reflejan en mejores resultados académicos, lo que sugiere la existencia de variables adicionales que condicionan el impacto de la tecnología en la educación.
\subsection{Objetivo del negocio}
Desarrollar un análisis basado en datos que permita identificar la relación entre el acceso a internet, las condiciones socioeconómicas y el desempeño educativo en los municipios colombianos, con el fin de generar insumos que apoyen la toma de decisiones del Ministerio de Educación Nacional en la priorización de inversiones y estrategias de intervención.
\subsection{Objetivos específicos}
Analizar la relación entre conectividad y resultados en Saber 11
Evaluar el impacto de la pobreza multidimensional en el desempeño académico
Identificar brechas territoriales en educación y acceso digital
Detectar patrones y agrupaciones de municipios con características similares
Establecer una base analítica para el desarrollo de modelos en la siguiente fase